\documentclass[a4paper,12pt]{report}

\usepackage[lmargin=2.00000cm,rmargin=2.0000cm,tmargin=5.5000cm,bmargin=2.500000cm,headheight=3.5cm]{geometry}        %Flexible and complete interface to document dimensions
\usepackage[utf8]{inputenc}
\usepackage[T1]{fontenc}
%\usepackage[latin1]{inputenc}
\usepackage{amsmath}
\usepackage{amsfonts}
\usepackage{amssymb}
\usepackage{lmodern}
\usepackage{float}
\usepackage{graphicx}

\usepackage[english,french]{babel}        %use for the below package `datetime'
\usepackage[babel=true]{csquotes}
\usepackage{datetime}        %Change format of `\today' with commands for current time
\renewcommand{\dateseparator}{-}
\newcommand{\headertoday}{\twodigit\day \dateseparator \twodigit\month  \dateseparator \the\year}

\usepackage{animate}
\usepackage{lastpage}
\usepackage{array}
\usepackage{lastpage}
\usepackage{multirow}
\usepackage{titling}
\usepackage{placeins}
\usepackage{media9}
\usepackage{eurosym}
\usepackage{pdflscape}
\usepackage{color}
\usepackage[table]{xcolor}
\definecolor{mon_bleu}{rgb}{0.137,0.466,0.741}
\definecolor{mon_vert}{rgb}{0.07843,0.4627,0.07843}
\definecolor{mon_rouge}{rgb}{0.62745,0.16078,0.27058}		
%\usepackage{subfigure}
\usepackage{subcaption}

\usepackage[ 
           hidelinks,
           colorlinks=true,
           linkcolor=blue,          % color of internal links (change box color with linkbordercolor)
    	   citecolor=blue,        % color of links to bibliography
    	   filecolor=blue,      % color of file links
    	   urlcolor=blue,           % color of external links
           pdfhighlight =/O]{hyperref}																		% dvipdfm package pour lien hypertextes dans pdf (pdfhighlight--> afficher la main dans le pdf, colorlinks--> colorlinks=true afficher les liens en couleurs)
  
\usepackage{blindtext}
\usepackage[final]{pdfpages}
\usepackage[french]{cleveref}																				% à placer après hyperref attention ne pas utiliser ":" dans les labels des équations avec French babel activé et cref

\usepackage{fancyhdr}
%%%%%%%%%%%%%%%%%%%%%%%%%%%%%%%%%%%%%%%%%%%%%%%%%%%%%%%%%%%%%%%%%%%%%%%%%%
%                         FORMAT PERSONNALISES                           %
% %%%%%%%%%%%%%%%%%%%%%%%%%%%%%%%%%%%%%%%%%%%%%%%%%%%%%%%%%%%%%%%%%%%%%%%%

\parindent=0pt        %leading space for paragraphs
\pagestyle{fancy}
\renewcommand{\arraystretch}{1.5}
\renewcommand{\headrulewidth}{0pt}
\fancyhead[CE,CO,LE,LO,RE,RO]{} %% clear out all headers
\fancyhead[C]{%
\begin{tabular}{|m{3.0cm}|m{10cm}|c@{}|}
\hline
\multirow{2}{*}{\includegraphics[scale=0.045]																			% LOGO LABO
{Logo-Symme.png}}& \centering \multirow{2}{*}{ \Large{\thetitle}}  &  \multirow{2}{*}{Page: \thepage ~/ \pageref{LastPage}}\\
& &  \\
\hline
%Etabli par: \newline \theauthor & \multicolumn{2}{l|}{Diffusion: interne}\\
%\hline
\multicolumn{2}{|l|}{Objet: \thetitleobject }& Date: \headertoday ~~\\ 
\hline
\end{tabular}
}
\renewcommand\footrulewidth{1pt}
\fancyfoot[L]{Document confidentiel} 																		% DOCUMENT CONFIDENTIEL
\fancyfoot[R]{Laboratoire SYMME}		 																	% NOM DU LABO

\newcolumntype{x}[1]{>{\centering\hspace{0pt}}p{#1}}
\setlength{\doublerulesep}{\arrayrulewidth} 


\usepackage{tikz}																							% pour les diagrammes
\usetikzlibrary{calc}
\usetikzlibrary{shapes.geometric,shapes.arrows,decorations.pathmorphing}
\usetikzlibrary{matrix,chains,scopes,positioning,shapes,arrows,fit}
\usetikzlibrary{calc,decorations.pathreplacing}
\usetikzlibrary{calc}
\usetikzlibrary{backgrounds,decorations.markings}

\setcounter{secnumdepth}{5}		%sous sous sous section



%%%%%%%%%%%%%%%%%%%%%%%%%%%%%%%%%%%%%%%%%%%%%%%%%%%%%%%%%%%%%%%%%%%%%%%%%%%%%%%%%%%%
%%-------------> PAGE DE GARDE INFO
%%%%%%%%%%%%%%%%%%%%%%%%%%%%%%%%%%%%%%%%%%%%%%%%%%%%%%%%%%%%%%%%%%%%%%%%%%%%%%%%%%%%

\author{Auteur}
\newcommand{\validator}{J. COLLOMB}
\title{Rapport d'avancement}
\selectlanguage{french}	
\date{\today}
\newcommand{\thetitleobject}{Comment intégrer des figures ?}
\setcounter{tocdepth}{6}
\setcounter{secnumdepth}{6}


%%%%%%%%%%%%%%%%%%%%%%%%%%%%%%%%%%%%%%%%%%%%%%%%%%%%%%%%%%%%%%%%%%%%%%%%%%%%%%%%%%%%
%%-------------> DEBUT DU DOCUMENT 
%%%%%%%%%%%%%%%%%%%%%%%%%%%%%%%%%%%%%%%%%%%%%%%%%%%%%%%%%%%%%%%%%%%%%%%%%%%%%%%%%%%%

\begin{document}

\graphicspath{{Figures/}}

%%%%%%%%%%%%%%%%%%%%%%%%%%%%%%%%%%%%%%%%%%%%%%%%%%%%%%%%%%%%%%%%%%%%%%%%%%%%%%%%%%%%
%%-------------> PAGE DE GARDE 
%%%%%%%%%%%%%%%%%%%%%%%%%%%%%%%%%%%%%%%%%%%%%%%%%%%%%%%%%%%%%%%%%%%%%%%%%%%%%%%%%%%%
\begin{titlepage}
    \centering
    \vspace*{3cm}
    {\bfseries\Large
        \thetitle\\
        \thetitleobject\\
        --- \\
        Document confidentiel\\  																			% LIGNE "CONFIDENTIEL"
        \vskip2cm
    }
    \includegraphics[scale=0.65]{logos.png}  																% LOGOS DES PARTENAIRES
    \vfill
    \begin{center}
    \begin{table}[b]
    \begin{tabular}{x{.225\linewidth}|| x{.225\linewidth}|| x{.225\linewidth} || x{.225\linewidth} }
    \hline \hline
    \textbf{Date} & \textbf{Révision} & \textbf{Rédigé par} & \textbf{Validé par} \tabularnewline
    \hline \hline
    \thedate & Rev. A  &  \theauthor &  \validator \tabularnewline
    \hline
    ~ & ~ & ~ & ~ \tabularnewline
    \hline
    ~ & ~ & ~ & ~ \tabularnewline
    \hline
    ~ & ~ & ~ & ~\tabularnewline         
    \hline \hline                                                                   
    \end{tabular}
    \end{table}
    \end{center}
    \vfill
    \vfill
    \begin{tikzpicture}[remember picture,overlay]
    \node (label) at (8cm,20cm){
        \includegraphics[width=2.5cm]{Logo-Symme.png} 													% LOGO LABO
      };
    
	\draw[very thick]
		([yshift=-25pt,xshift=25pt]current page.north west)--
		([yshift=-25pt,xshift=-25pt]current page.north east)--
		([yshift=25pt,xshift=-25pt]current page.south east)--
		([yshift=25pt,xshift=25pt]current page.south west)--cycle;
    \end{tikzpicture}
\end{titlepage}


%\begin{titlepage}
%    \centering
%    
%    \vspace*{3cm}
%    {\bfseries\Large
%        \thetitle\\
%        \thetitleobject\\
%        \vskip2cm
%    }    
%    \vfill
%    \begin{center}
%    \begin{table}[b]
%    \begin{tabular}{x{.225\linewidth}|| x{.225\linewidth}|| x{.225\linewidth} || x{.225\linewidth} }
%    \hline \hline
%    \textbf{Date} & \textbf{R�vision} & \textbf{R�dig� par} & \textbf{Valid� par} \tabularnewline
%    \hline \hline
%    \thedate & Rev. A  &  \theauthor &  \validator \tabularnewline
%    \hline
%    ~ & ~ & ~ & ~ \tabularnewline
%    \hline
%    ~ & ~ & ~ & ~ \tabularnewline
%    \hline
%    ~ & ~ & ~ & ~\tabularnewline         
%    \hline \hline                                                                   
%    \end{tabular}
%    \end{table}
%    \end{center}
%    \vfill
%    \vfill
%    \begin{tikzpicture}[overlay,remember picture]
%    \node (label) at (6.7cm,19.5cm){
%        \includegraphics[width=4cm]{../Figures/logo_CT1.pdf} % also works with logo.pdf
%      };
%    
%    \draw [line width=1pt,rounded corners=7pt]
%        ($ (current page.north west) + (1.5cm,-1.5cm) $)
%        rectangle
%        ($ (current page.south east) + (-1.5cm,1.5cm) $);
%        
%%	\draw[very thick]
%%		([yshift=-25pt,xshift=25pt]current page.north west)--
%%		([yshift=-25pt,xshift=-25pt]current page.north east)--
%%		([yshift=25pt,xshift=-25pt]current page.south east)--
%%		([yshift=25pt,xshift=25pt]current page.south west)--cycle;
%    \end{tikzpicture}
%\end{titlepage}


%%%%%%%%%%%%%%%%%%%%%%%%%%%%%%%%%%%%%%%%%%%%%%%%%%%%%%%%%%%%%%%%%%%%%%%%%%%%%%%%%%%%
%%-------------> SOMMAIRE
%%%%%%%%%%%%%%%%%%%%%%%%%%%%%%%%%%%%%%%%%%%%%%%%%%%%%%%%%%%%%%%%%%%%%%%%%%%%%%%%%%%%

\renewcommand\contentsname{Sommaire}
\setcounter{chapter}{1}
\tableofcontents
%\listoffigures
%\listoftables


%%%%%%%%%%%%%%%%%%%%%%%%%%%%%%%%%%%%%%%%%%%%%%%%%%%%%%%%%%%%%%%%%%%%%%%%%%%%%%%%%%%%
%%-------------> CORPS DOCUMENT
%%%%%%%%%%%%%%%%%%%%%%%%%%%%%%%%%%%%%%%%%%%%%%%%%%%%%%%%%%%%%%%%%%%%%%%%%%%%%%%%%%%%

%-----------------------------------------------------------------------------------
%-----------------------------------------------------------------------------------
\newpage

\section{Vous avez dit figures ?}
\subsection{Image matricielle et image vectorielle}

\textit{\textbf{Une image matricielle}, ou « carte de points » (de l'anglais bitmap), est une image constituée d'une matrice de points colorés. C'est-à-dire, constituée d'un tableau, d'une grille, où chaque case possède une couleur qui lui est propre et est considérée comme un point. Il s'agit donc d'une juxtaposition de points de couleurs formant, dans leur ensemble, une image.} \\
Source : \href{https://fr.wikipedia.org/wiki/Image_matricielle}{Wikipédia - Image matricielle} \\

\textit{\textbf{Une image vectorielle} (ou image en mode trait), en informatique, est une image numérique composée d'objets géométriques individuels, des primitives géométriques (segments de droite, arcs de cercle, courbes de Bézier, polygones, etc.), définis chacun par différents attributs (forme, position, couleur, remplissage, visibilité, etc.) et auxquels on peut appliquer différentes transformations (homothéties, rotations, écrasement, mise à l'échelle, extrusion, inclinaison, effet miroir, dégradé de formes, morphage, symétrie, translation, interpolation, coniques ou bien les formes de révolution).}\\
Source : \href{https://fr.wikipedia.org/wiki/Image_vectorielle}{Wikipédia - Image vectorielle}


\subsection{Comment créer ses figures ?}
Il est courant pour les doctorants de devoir exploiter des résultats issus d'expérimentations :
\begin{itemize}
\item relevés de température;
\item mesures de déplacement;
\item corrélation d'images;
\item relevé GPS;
\item \ldots
\end{itemize}

\ \\

Le traitement des résultats peut s'effectuer par divers logiciels :
\begin{itemize}
\item Excel;
\item Matlab;
\item Python;
\item Orange;
\item \ldots
\end{itemize}

\ \\

\href{https://www.anaconda.com/download/}{Python} est un logiciel de programmation libre avec lequel il est possible d'exploiter divers résultats, d'automatiser des actions répétitives et d'afficher les données sous forme de Figure (\href{https://github.com/JeanCollomb/Python_plot}{Quelques exemples simples}). \\

Python permet l'export d'images matricielles (\verb|.jpg, .png|), mais également d'images vectorielles (\verb|.svg, .eps|).


\subsection{Créer des images sous Powerpoint}
Powerpoint est un outil simple et efficace pour la création de figures (algorithmes, annotation de courbes, \ldots). \\

Powerpoint permet l'export des figures ainsi générées au format \verb|.emf| (\href{http://saf.bio.caltech.edu/PPT_G_P_I/}{démarche}). Ces images peuvent par la suite être ouvertes par des logiciels libres tel que \href{https://inkscape.org/fr/}{Inkscape} puis exportée dans un format d'image vectorielle ou mieux, au format \verb|.pdf_tex|. Ce format permet d'encapsuler les images dans un pdf et de générer les annotations, géométries à l'aide d'un fichier Tex modifiable. \\

\underline{\textbf{Suggestion :}} \\
Pour les toutes les figures :
\begin{enumerate}
\item Préférer la génération d'images vectorielles (\verb|.eps|, \verb|.svg|);
\item (Option) Ajouter des compléments via Powerpoint et exporter au format \verb|.emf|;
\item Convertir les images en fichiers \verb|.pdf_tex| à l'aide du logiciel \href{https://inkscape.org/fr/}{Inkscape};
\item Intégrer uniquement des images \verb|.pdf_tex| dans le document \LaTeX (minimiser les formats différents pour les images permet de minimiser les sources d'erreurs).
\end{enumerate}



\section{Intégrer des figures sous \LaTeX}
\subsection{Figure simple}
Pour l'intégration des sous-figures, il est nécessaire de placer en début de document le package subcaption (le package subfig étant obsolète désormais). \\
La commande est donc : \verb|\usepackage{subcaption}|.

\subsection{Figure simple}
Le code \LaTeX permettant l'ajout d'une figure simple au format \verb|.pdf_tex| est le suivant : \\

\verb|\begin{figure}[hbtp]	| \\
\verb|	\centering 	% Figure centrée| \\
\verb|	\def\svgwidth{0.8\columnwidth} % Figure faisant 80% de la largeur de page | \\
\verb|	\fontsize{10pt}{10pt}\selectfont\input{Figures/Figure_1.pdf_tex} | \\
\verb|	\caption{Exemple de figure simple} | \\
\verb|	\label{figure_simple} | \\
\verb|\end{figure} | \\

\subsection{Sous-figures}
Le code \LaTeX permettant l'ajout d'une figure multiple au format \verb|.pdf_tex| est le suivant : \\

\verb|\begin{figure}[hbtp]| \\
\verb|	\centering| \\
\verb|	\begin{subfigure}[b]{0.8\textwidth}| \\
\verb|		\centering| \\
\verb|		\def\svgwidth{\columnwidth}| \\
\verb|		\fontsize{10pt}{10pt}\selectfont\input{Figures/Figure_2.pdf_tex}| \\
\verb|		\caption{Exemple subplot} | \\
\verb|		\label{figure_subplot}| \\
\verb|	\end{subfigure}| \\
\verb|	\qquad| \\
\verb|	\begin{subfigure}[b]{0.7\textwidth}| \\
\verb|		\centering| \\
\verb|		\def\svgwidth{\columnwidth}| \\
\verb|		\fontsize{10pt}{10pt}\selectfont\input{Figures/Figure_3.pdf_tex}| \\
\verb|		\caption{Exemple diagramme polaire} | \\
\verb|		\label{figure_polaire}| \\
\verb|	\end{subfigure}| \\
\verb|	\caption{Exemple de Figures multiple}| \\
\verb|	\label{figure_multiple}| \\
\verb|\end{figure}| \\



\newpage
\section{Exemples}
La Figure~\ref{figure_simple} est un exemple d'intégration d'une figure simple. \\
Les images sont au format vectoriel. Ça donne quoi un zoom dessus ?

\begin{figure}[hbtp]
	\centering
	\def\svgwidth{1\columnwidth}
	\fontsize{10pt}{10pt}\selectfont\input{Figures/Figure_1.pdf_tex}
	\caption{Exemple de figure simple}
	\label{figure_simple}
\end{figure}



La Figure~\ref{figure_multiple} est un exemple d'intégration d'une figure multiple. \\
Il est bien entendu possible de faire référence à la Figure~\ref{figure_subplot} ou à la Figure~\ref{figure_polaire}.

\begin{figure}[hbtp]
	\centering
	\begin{subfigure}[b]{0.8\textwidth}
		\centering
		\def\svgwidth{\columnwidth}
		\fontsize{10pt}{10pt}\selectfont\input{Figures/Figure_2.pdf_tex}
		\caption{Exemple subplot} 
		\label{figure_subplot}
	\end{subfigure}
	\qquad
	\begin{subfigure}[b]{0.7\textwidth}
		\centering
		\def\svgwidth{\columnwidth}
		\fontsize{10pt}{10pt}\selectfont\input{Figures/Figure_3.pdf_tex}
		\caption{Exemple diagramme polaire} 
		\label{figure_polaire}
	\end{subfigure}
	\caption{Exemple de Figures multiple} 
	\label{figure_multiple}
\end{figure}



%%%%%%%%%%%%%%%%%%%%%%%%%%%%%%%%%%%%%%%%%%%%%%%%%%%%%%%%%%%%%%%%%%%%%%%%%%%%%%%%%%%%
%%-------------> FIN DU DOCUMENT
%%%%%%%%%%%%%%%%%%%%%%%%%%%%%%%%%%%%%%%%%%%%%%%%%%%%%%%%%%%%%%%%%%%%%%%%%%%%%%%%%%%%

\end{document}